% ---- pg-380 （书页 368） ----

顺序, 列表表示每个特征值对样本总方差的贡献百分比.

(c) 已知(b)的结果, 确定重要样本主成分的个数, 若用单一的反应指标成分来综合放射疗法数据是否可行? 请说明.

(d) 列表写出每一个你想保留的主成分与原始数据间的相关系数. 如果可能, 解释这些成分.

\noindent 8.14\quad 用例 5.2 的汗液数据的样本协方差矩阵作主成分分析, 作每个重要主成分的 Q-Q 散点图, 有可疑的观察值吗? 试说明.

\noindent 8.15\quad 在例 8.6 中讨论过的生育后体重的 4 个样本标准差为
\[
\sqrt{s_{11}}=32.990\,9,\quad \sqrt{s_{22}}=33.591\,8,\quad \sqrt{s_{33}}=36.553\,4,\quad \sqrt{s_{44}}=37.351\,7
\]
用这些数据和例 8.6 给定的相关系数, 构造样本协方差矩阵 $\mathbf{S}$, 用 $\mathbf{S}$ 作主成分分析.

\noindent 8.16\quad 20 世纪 90 年代, 在五年的时间里, 以威斯康星的 28 个湖上的渔民作为年度样本, 要求他们报告捕鱼花费的时间及捕到的各种类型的鱼的数量, 然后把他们反馈的信息转化成对下面几种鱼的每小时的捕捉率,
\[
\begin{array}{lll}
x_1=\text{蓝鳃鱼} & x_2=\text{黑刺盖太阳鱼} & x_3=\text{小嘴鲈鱼}\\
x_4=\text{大嘴鲈鱼} & x_5=\text{突眼鱼} & x_6=\text{北方狗鱼}
\end{array}
\]
估计的相关系数矩阵为
\[
\mathbf{R}=\begin{bmatrix}
1&0.491\,9&0.263\,6&0.465\,3&-0.227\,7&0.065\,2\\
0.491\,9&1&0.312\,7&0.350\,6&-0.191\,7&0.204\,5\\
0.263\,5&0.312\,7&1&0.410\,8&0.064\,7&0.249\,3\\
0.465\,3&0.350\,6&0.410\,8&1&-0.224\,9&0.229\,3\\
-0.227\,7&-0.191\,7&0.064\,7&-0.224\,9&1&-0.214\,4\\
0.065\,2&0.204\,5&0.249\,3&0.229\,3&-0.214\,4&1
\end{bmatrix}
\]
它是基于一个容量为 120 的样本(有少量数值丢失了).

同一个渔夫捕到的鱼都放在一起, 因此这些数据可以提供关于鱼群的一些迹象. 前四种鱼属于棘臀太阳鱼科, 这是数量最丰富的科. 突眼鱼则是最受欢迎的食用鱼.

(a) 论述棘臀太阳鱼科内 $x_1$ 至 $x_4$ 的相关系数模型, 突眼鱼与其他几种鱼是同一鱼群的吗?

(b) 仅用 $x_1$ 至 $x_4$ 作主成分分析, 解释你的结果.

(c) 用全部 6 个变量进行主成分分析. 解释得到的结果.

\noindent 8.17\quad 用\hyperref[tab:1.8]{表 1.8} 中的骨头所含矿物质的数据, 作 $\mathbf{S}$ 的主成分分析.

\noindent 8.18\quad 国家和地区的女子田径纪录数据列于\hyperref[tab:1.9]{表 1.9}.

(a) 求这些数据的样本相关系数矩阵 $\mathbf{R}$, 并求它的特征值和特征向量.

(b) 求标准化变量的前两个主成分. 列表表示标准化变量的相关系数, 并给出主成分及由这两个主成分解释的(标准化)样本总方差累积百分比.

(c) 解释(b)中求得的两个主成分(注意, 第一成分实质上是规范化的单位向量, 并且可以测量给定国家和地区的运动员的水平. 第二成分测量国家和地区在不同赛跑距离中的相对实力).

(d) 把国家和地区按它们在第一主成分的得分排列, 这种排列与你最初对不同国家和地区运动水平的看法相一致吗?
